\fancyhead[L]{\songti\zihao{-5}\cquHeader}
\fancyhead[R]{\songti\zihao{-5}摘要}
\addcontentsline{toc}{section}{摘要}
\section*{摘\quad 要}
\zihao{-4}

形式化数学通过交互式定理证明器把数学命题和证明转化为可由计算机检查的严格表达，是提高复杂证明可靠性的重要方法。随着大语言模型和证明智能体的发展，人工智能开始参与数学证明生成、Lean 代码补全、定理检索和错误修复等环节。本文围绕人工智能辅助 Lean 4 形式化证明展开研究，并以 Mordell 方程整数解问题作为具体案例，考察传统数论证明在 Lean 4 与 Mathlib4 环境中的形式化过程。

本文首先梳理 AI 数学推理和形式化证明系统的发展现状，然后围绕 Mordell 方程
\[
y^2=x^3+d
\]
建立 Lean 4 形式化项目。证明主线从二次整数环 $\mathbb{Z}[\sqrt d]$ 中的元素分解出发，将方程转化为主理想分解，再利用 Dedekind 整环结构、共轭主理想互素性、类群下降和有限类群计算，得到 Mordell 基本结论，并将其应用于 $d=-1,-2,-5,-6,-13$ 等具体参数。最终项目包含 12 个 Lean 源文件，约 2800 行代码，完成了基础二次整数环接口、Dedekind 条件、类群计算以及最终整数解定理的形式化；项目通过 \texttt{lake build} 验证，未使用 \texttt{sorry}、\texttt{admit} 或 \texttt{axiom} 占位。

在此基础上，本文进一步对开发过程中的人机交互记录进行任务级评价。本文将 47 个会话、413 条消息划分为 76 个任务，并参照评价框架标注 Full、Partial、Incorrect 和 Pending 等完成结果。统计显示，AI 智能体在本项目中的 Full 比例为 77.6\%，Full 与 Partial 合计的有效贡献比例为 92.1\%。实验表明，当前 AI 能够有效提供证明路线、局部引理、代码草稿和编译错误修复建议，但在复杂代数数论形式化任务中仍难以脱离人类设计的证明结构和 Lean 编译反馈独立完成完整证明。本文的实践说明，人类给出数学路线、AI 辅助证明工程、Lean 负责最终核验的协作模式，已经能够支持具有一定复杂度的形式化数学项目。

~\\
\hspace*{2em}{\heiti\zihao{-4}关键词}：人工智能；Lean 4；形式化证明；Mordell 方程；

\newpage
\fancyhead[L]{\songti\zihao{-5}\cquHeader}
\fancyhead[R]{\zihao{-5}ABSTRACT}
\addcontentsline{toc}{section}{ABSTRACT}
\titleformat{\section}[block]{\centering\bfseries\fontspec{Times New Roman}\fontsize{16pt}{20pt}\selectfont}{\thesection}{1em}{}[]
\section*{ABSTRACT}
\zihao{-4}
\setmainfont{Times New Roman}

Formal mathematics transforms mathematical statements and proofs into rigorous expressions that can be checked by computers through interactive theorem provers. It is an important method for improving the reliability of complex proofs. With the development of large language models and proof agents, artificial intelligence has begun to participate in proof generation, Lean code completion, theorem retrieval, and error correction. This thesis studies AI-assisted formal proof development in Lean 4, using the problem of integral solutions of Mordell equations as a concrete case. It examines how a traditional number-theoretic proof can be formalized in the Lean 4 and Mathlib4 environment.

This thesis first reviews the recent development of AI-based mathematical reasoning and formal proof systems. It then constructs a Lean 4 formalization project for the Mordell equation
\[
y^2=x^3+d.
\]
The main proof starts from a factorization in the quadratic integer ring $\mathbb{Z}[\sqrt d]$. The equation is transformed into a factorization of principal ideals, and the argument then uses Dedekind domain structures, the coprimality of conjugate principal ideals, class group descent, and finite class group computations. These steps lead to the basic Mordell result and are further applied to the concrete parameters $d=-1,-2,-5,-6,-13$. The final project consists of 12 Lean source files and about 2800 lines of code. It formalizes the basic interface of quadratic integer rings, Dedekind conditions, class group computations, and the final theorems on integral solutions. The project is verified by \texttt{lake build} and does not use \texttt{sorry}, \texttt{admit}, or \texttt{axiom} as placeholders.

On this basis, the thesis also evaluates the human-agent interaction records from the development process at the task level. The records contain 47 sessions and 413 messages, which are grouped into 76 tasks and annotated using the Full, Partial, Incorrect, and Pending categories of the evaluation framework. The resulting Full rate is 77.6\%, and the combined effective contribution rate of Full and Partial tasks is 92.1\%. The experiments show that current AI systems can effectively provide proof outlines, local lemmas, code drafts, and suggestions for fixing compilation errors. However, in complex formalization tasks involving algebraic number theory, they still cannot complete a full proof independently without a human-designed proof structure and feedback from the Lean compiler. This case study shows that the collaborative workflow in which humans provide the mathematical strategy, AI assists with proof engineering, and Lean performs the final verification is already capable of supporting formal mathematics projects of moderate complexity.

~\\
\hspace*{2em}\textbf{Key words}: artificial intelligence; Lean 4; formal proof; Mordell equation

\setmainfont{Times New Roman}
\titleformat{\section}{\heiti\zihao{3}\centering}{\thesection}{0.5em}{}[]
